\documentclass[10pt,a4paper]{article}
\usepackage{cmap} % ToUnicode-таблицы: без них кириллица в PDF не ищется и не копируется
\usepackage[utf8]{inputenc}
\usepackage[T1,T2A]{fontenc}
\usepackage[russian]{babel}
\usepackage{amsmath}
\usepackage[margin=1in]{geometry}
\usepackage{csquotes}

% Копия макросов из lecture_preamble.tex: эта памятка не подключает преамбулу слайдов.
% Короткое вертикальное многоточие: размеры в em, поэтому оно уменьшается вместе со шрифтом
% (у \vdots они фиксированы в pt).
\DeclareRobustCommand{\ndvdots}{%
    \mathord{\vbox{\baselineskip=.3em\lineskiplimit=0pt\kern.15em\hbox{.}\hbox{.}\hbox{.}}}}
% Подвывод в схеме правила, нарисованный как во fitch: вертикальная черта и флажок под допущением.
% Выравнивание по нижней строке, чтобы в \lor E заключения подвыводов стояли на одной линии с A \lor B.
\newcommand{\ndsubproof}[2]{{%
    \setlength{\arrayrulewidth}{.2mm}%
    \begin{array}[b]{@{}|@{\,}l@{\,}}#1\\\hline\ndvdots\\#2\end{array}}}
\newcommand{\RAA}{\mathit{RAA}}

\begin{document}
\begin{itemize}
    \item Для $\land$:
    \[
    \frac{A \qquad B}{A \wedge B}\ \wedge I
    \qquad \qquad
    \frac{A \wedge B}{A}\ \wedge{E}
    \qquad
    \frac{A \wedge B}{B}\ \wedge{E}
    \]
    \item Для $\to$:
    \[
    \frac{\ndsubproof{A}{B}}{A \to B}\ \to I
    \qquad \qquad
    \frac{A \to B \quad A}{B}\ \to E
    \]
    \item Для $\neg$ и $\bot$:
    \[
    \frac{\ndsubproof{A}{\bot}}{\neg A}\ \neg I
    \qquad \qquad
    \frac{\neg A \quad A}{\bot}\ \neg E/\bot I
    \qquad \qquad
    \frac{\bot}{A}\ \bot{E}
    \]
    \item Для $\lor$:
    \[
    \frac{A}{A \lor B}\ \lor{I}
    \qquad
    \frac{B}{A \lor B}\ \lor{I}
    \qquad \qquad
    \frac{A \lor B \quad \ndsubproof{A}{C} \quad \ndsubproof{B}{C}}{C}\ \lor E
    \]
    \item Остальные:
    \[
    \frac{A}{A}\ R
    \qquad \qquad
    \frac{\ndsubproof{\neg A}{\bot}}{A}\ \RAA
    \qquad \qquad
    \frac{\ndsubproof{\neg A}{B}}{A \lor B}\ \lor I'
    \qquad
    \frac{\ndsubproof{\neg B}{A}}{A \lor B}\ \lor I'
    \]
    \item $I$ "--- введение, $E$ "--- исключение. Например, $\land I$ "--- \enquote{введение конъюнкции} (также В$\land$).
    Правило введения показывает, как доказать формулу с данной главной связкой, правило исключения "--- как её использовать. Так, у $\land I$ формула $A\land B$ "--- заключение, а у $\land E$ "--- посылка.
\end{itemize}
\end{document}
